La Loi Générale des Masses est formulée comme une transformation généalogique et
spectrale, et non comme une collection de nombres. Son domaine effectif est la chaîne
\[
\text{APP}\longrightarrow\text{TRIT}\longrightarrow\text{TPK}
\longrightarrow\Gamma^{\rm surv}\longrightarrow\Gamma^{\rm obs}
\longrightarrow L_{108}\longrightarrow\sigma\in\mathbb Z^5
\longrightarrow\chi_{\rm full}\longrightarrow\widehat M,
\]
suivie, lorsque la représentation physique le permet, de la fonctionnelle qui
extrait une coordonnée scalaire de la droite de masse. Chaque flèche conserve son
domaine, son orientation, sa mémoire et son statut ; aucune valeur externe ne choisit
rétroactivement un parcours, une signature ou un coefficient harmonique.

L'atlas fini contient 81 cellules, 56 familles de parcours, 13 multisections et
324 parcours orientés. Le point central électronique est l'unique fibre ponctuelle
équivariante. Les autres classes élémentaires se réalisent naturellement comme
fibres, orbites ou multisections, si bien que leur résultat primaire peut être un
opérateur ou un spectre avant d'être une masse scalaire. Cette différence n'est pas une
lacune du domaine, mais une propriété de la géométrie interne.

Dans la fibre centrale, le sélecteur de base et le lecteur bidirectionnel produisent
\[
\mathfrak e_0
=\frac{\sqrt3}{4}
 \left[\exp\!\left(\frac{\varphi}{\pi^2}\right)
       -22\alpha_{\rm HMT}^3\right](1+15\Delta_4),
\qquad
\kappa_e=80-54\alpha_{\rm HMT}+6\Delta_4,
\]
\[
\mathfrak e_e^\beta=\mathfrak e_0R_{\rm act}^{\kappa_e},
\qquad
\varsigma_{u_E}(\mathfrak e_e^\beta)
=0.510998950690000808608\ldots\ \mathrm{MeV}.
\]
L'échelle absolue d'action demeure dans la branche
\(\Sigma_H=\varphi\alpha_{\rm HMT}^{16}\), d'ordre décimal \(-34\), et se
couple à l'horloge par \(\hbar_{\rm int}\omega\) ; dans les quotients de
masses, toute section commune se simplifie et seuls subsistent les caractères, les
lecteurs et les raffinements relatifs correctement typés.

La hiérarchie harmonique ne consiste pas en huit corrections isolées. C'est le semi-groupe
global
\[
\langle90,120\rangle
=\{90,120\}\cup\{30r:r\ge6\},
\qquad
c_n=q_-^n+q_+^n,
\quad s_n=q_-^n-q_+^n,
\]
agissant sur un unique opérateur bicouche. Les niveaux visibles sont des témoins de
cette hiérarchie ; sa complétude représentative ne remplace pas l'application physique qui
doit transporter parcours, retenue, mémoire et enroulement vers les coefficients de chaque
état.

Les réalisations physiques typées, les dix-sept évaluations conditionnelles, les
deux sections nulles et les secteurs neutriniques, composés, résonants,
topologiques et virtuels sont ainsi réunis sans aplanir leurs codomaines.
L'inventaire ultérieur de 471 masses et de 384 largeurs conserve la confrontation
état par état, mais ne se confond pas avec un dénombrement de prédictions. Une
comparaison quantitative n'acquiert une force physique qu'après déclaration de
l'observable, de l'unité, du schéma, de l'échelle, de l'incertitude, de la covariance et de la règle de
réalisation.

La conclusion structurelle est précise : la masse est la mesure spectrale d'une
persistance sous couplage. La particule est la section persistante ; le parcours,
son histoire observable ; la signature, sa lecture entière ; le caractère, son action
adimensionnelle ; les tours, sa résolution harmonique ; et la carte dimensionnelle, la
réalisation de cette généalogie en unités physiques. L'atlas et les inventaires
précédents permettent de reproduire chaque affirmation sans réduire la théorie à un
tableau abrégé ni transformer une interface ouverte en théorème établi.
